\documentclass[
  man,
  floatsintext,
  longtable,
  nolmodern,
  notxfonts,
  notimes,
  colorlinks=true,linkcolor=blue,citecolor=blue,urlcolor=blue]{apa7}

\usepackage{amsmath}
\usepackage{amssymb}
\usepackage{setspace}



\RequirePackage{longtable}
\RequirePackage{threeparttablex}

\usepackage{framed}


\makeatletter
\renewcommand{\paragraph}{\@startsection{paragraph}{4}{\parindent}%
	{0\baselineskip \@plus 0.2ex \@minus 0.2ex}%
	{-.5em}%
	{\normalfont\normalsize\bfseries\typesectitle}}

\renewcommand{\subparagraph}[1]{\@startsection{subparagraph}{5}{0.5em}%
	{0\baselineskip \@plus 0.2ex \@minus 0.2ex}%
	{-\z@\relax}%
	{\normalfont\normalsize\bfseries\itshape\hspace{\parindent}{#1}\textit{\addperi}}{\relax}}
\makeatother




\usepackage{longtable, booktabs, multirow, multicol, colortbl, hhline, caption, array, float, xpatch}
\usepackage{subcaption}


\renewcommand\thesubfigure{\Alph{subfigure}}
\setcounter{topnumber}{2}
\setcounter{bottomnumber}{2}
\setcounter{totalnumber}{4}
\renewcommand{\topfraction}{0.85}
\renewcommand{\bottomfraction}{0.85}
\renewcommand{\textfraction}{0.15}
\renewcommand{\floatpagefraction}{0.7}

\usepackage{tcolorbox}
\tcbuselibrary{listings,theorems, breakable, skins}
\usepackage{fontawesome5}

\definecolor{quarto-callout-color}{HTML}{909090}
\definecolor{quarto-callout-note-color}{HTML}{0758E5}
\definecolor{quarto-callout-important-color}{HTML}{CC1914}
\definecolor{quarto-callout-warning-color}{HTML}{EB9113}
\definecolor{quarto-callout-tip-color}{HTML}{00A047}
\definecolor{quarto-callout-caution-color}{HTML}{FC5300}
\definecolor{quarto-callout-color-frame}{HTML}{ACACAC}
\definecolor{quarto-callout-note-color-frame}{HTML}{4582EC}
\definecolor{quarto-callout-important-color-frame}{HTML}{D9534F}
\definecolor{quarto-callout-warning-color-frame}{HTML}{F0AD4E}
\definecolor{quarto-callout-tip-color-frame}{HTML}{02B875}
\definecolor{quarto-callout-caution-color-frame}{HTML}{FD7E14}

%\newlength\Oldarrayrulewidth
%\newlength\Oldtabcolsep


\usepackage{hyperref}
\usepackage[protrusion=true]{microtype}



\providecommand{\tightlist}{%
  \setlength{\itemsep}{0pt}\setlength{\parskip}{0pt}}
\usepackage{longtable,booktabs,array}
\usepackage{calc} % for calculating minipage widths
% Correct order of tables after \paragraph or \subparagraph
\usepackage{etoolbox}
\makeatletter
\patchcmd\longtable{\par}{\if@noskipsec\mbox{}\fi\par}{}{}
\makeatother
% Allow footnotes in longtable head/foot
\IfFileExists{footnotehyper.sty}{\usepackage{footnotehyper}}{\usepackage{footnote}}
\makesavenoteenv{longtable}

\usepackage{graphicx}
\makeatletter
\newsavebox\pandoc@box
\newcommand*\pandocbounded[1]{% scales image to fit in text height/width
  \sbox\pandoc@box{#1}%
  \Gscale@div\@tempa{\textheight}{\dimexpr\ht\pandoc@box+\dp\pandoc@box\relax}%
  \Gscale@div\@tempb{\linewidth}{\wd\pandoc@box}%
  \ifdim\@tempb\p@<\@tempa\p@\let\@tempa\@tempb\fi% select the smaller of both
  \ifdim\@tempa\p@<\p@\scalebox{\@tempa}{\usebox\pandoc@box}%
  \else\usebox{\pandoc@box}%
  \fi%
}
% Set default figure placement to htbp
\def\fps@figure{htbp}
\makeatother







\usepackage{fontspec} 

\defaultfontfeatures{Scale=MatchLowercase}
\defaultfontfeatures[\rmfamily]{Ligatures=TeX,Scale=1}

  \setmainfont[,RawFeature={fallback=mainfontfallback}]{Times New Roman}




\title{Practice Effects Confound the Emotion--Performance Link in
Self-Regulated Learning: A Reanalysis of Lin et al.~(2026)}


\shorttitle{PRACTICE EFFECTS CONFOUND EMOTION--PERFORMANCE}


\usepackage{etoolbox}






\author{Chi Zhang}



\affiliation{
{Manchester Institute of Education, The University of Manchester}}




\leftheader{Zhang}



\abstract{Lin et al.~(2026) proposed a computational model of the
dynamic interplay between goal setting, performance, and emotions in
self-regulated learning. In a repetitive math task (Study 1), they
reported that positive emotions predicted lower subsequent performance,
interpreting this via the coasting hypothesis. We argue that this
finding may instead reflect a temporal confound: in a repetitive task,
performance improves due to practice while emotional engagement
declines, producing a spurious negative association when time trends are
unmodeled. We reanalyze the Study 1 data by introducing a log(trial)
learning curve into the best-fitting model. The negative
emotion-to-performance effect is no longer credible in this extended
model; instead, a strong practice effect emerges. This confound-based
interpretation also parsimoniously accounts for why the
emotion-to-performance pathway reversed direction in the original Study
2, where task variety and high stakes attenuate practice-driven time
trends. Methodologically, we show that the WAIC criterion used in the
original study exhibits poor diagnostic reliability and is theoretically
biased for autoregressive panel models. We implement exact
leave-future-out cross-validation as an unbiased alternative, under
which the original model selection conclusions are no longer
supported. }

\keywords{coasting hypothesis, computational modeling, leave-future-out
cross-validation, model comparison, practice effects, self-regulated
learning}

\authornote{\par{\addORCIDlink{Chi Zhang}{0000-0002-9535-6186}} 
\par{ }
\par{   The author states that there is no conflict of interest.  The
author would like to acknowledge the assistance given by Research IT,
and the use of The HPC Pool funded by the Research Lifecycle Programme
at The University of Manchester.  }
\par{Correspondence concerning this article should be addressed to Chi
Zhang, Manchester Institute of Education, The University of
Manchester, Oxford Road, Manchester, England M13
9PL, UK, Email: \href{mailto:chi.zhang-3@manchester.ac.uk}{chi.zhang-3@manchester.ac.uk}}
}

\makeatletter
\let\endoldlt\endlongtable
\def\endlongtable{
\hline
\endoldlt
}
\makeatother

\urlstyle{same}



\makeatletter
\@ifundefined{cslhangindent}{\newlength{\cslhangindent}}{}
\makeatother
\makeatletter
\@ifpackageloaded{caption}{}{\usepackage{caption}}
\AtBeginDocument{%
\ifdefined\contentsname
  \renewcommand*\contentsname{Table of contents}
\else
  \newcommand\contentsname{Table of contents}
\fi
\ifdefined\listfigurename
  \renewcommand*\listfigurename{List of Figures}
\else
  \newcommand\listfigurename{List of Figures}
\fi
\ifdefined\listtablename
  \renewcommand*\listtablename{List of Tables}
\else
  \newcommand\listtablename{List of Tables}
\fi
\ifdefined\figurename
  \renewcommand*\figurename{Figure}
\else
  \newcommand\figurename{Figure}
\fi
\ifdefined\tablename
  \renewcommand*\tablename{Table}
\else
  \newcommand\tablename{Table}
\fi
}
\@ifpackageloaded{float}{}{\usepackage{float}}
\floatstyle{ruled}
\@ifundefined{c@chapter}{\newfloat{codelisting}{h}{lop}}{\newfloat{codelisting}{h}{lop}[chapter]}
\floatname{codelisting}{Listing}
\newcommand*\listoflistings{\listof{codelisting}{List of Listings}}
\makeatother
\makeatletter
\makeatother
\makeatletter
\@ifpackageloaded{caption}{}{\usepackage{caption}}
\@ifpackageloaded{subcaption}{}{\usepackage{subcaption}}
\makeatother

% From https://tex.stackexchange.com/a/645996/211326
%%% apa7 doesn't want to add appendix section titles in the toc
%%% let's make it do it
\makeatletter
\xpatchcmd{\appendix}
  {\par}
  {\addcontentsline{toc}{section}{\@currentlabelname}\par}
  {}{}
\makeatother

%% Disable longtable counter
%% https://tex.stackexchange.com/a/248395/211326

\usepackage{etoolbox}

\makeatletter
\patchcmd{\LT@caption}
  {\bgroup}
  {\bgroup\global\LTpatch@captiontrue}
  {}{}
\patchcmd{\longtable}
  {\par}
  {\par\global\LTpatch@captionfalse}
  {}{}
\apptocmd{\endlongtable}
  {\ifLTpatch@caption\else\addtocounter{table}{-1}\fi}
  {}{}
\newif\ifLTpatch@caption
\makeatother

\begin{document}

\maketitle




\setlength\LTleft{0pt}



\ifdefined\Shaded\renewenvironment{Shaded}{\begin{snugshade}\begin{singlespace}}{\end{singlespace}\end{snugshade}}\fi

\setlength{\cslhangindent}{0.5in}

This is my first paragraph. Any section headings in the introduction
should be in levels 2--5.

\subsection{Section in Introduction}\label{section-in-introduction}

This is a reference.

\section{Method}\label{method}

\subsection{Participants}\label{participants}

\subsection{Measures}\label{measures}

\subsection{Procedure}\label{procedure}

\section{Results}\label{results}

\section{Discussion}\label{discussion}

\subsection{Model Comparison and Methodological
Considerations}\label{model-comparison-and-methodological-considerations}

We evaluated three models using multiple comparison criteria with
increasing methodological rigor. The original study employed WAIC for
model selection without reporting diagnostic checks. Our reanalysis
revealed that all three models exhibited substantial proportions of
unreliable pointwise WAIC estimates (\textgreater16\% p\_waic
\textgreater{} 0.4), calling into question the reliability of WAIC-based
model comparison in this context. We therefore additionally computed
PSIS-LOO, which, while more robust than WAIC (Vehtari et al., 2017), is
theoretically inappropriate for autoregressive models because leaving
out a single observation still allows future data to inform predictions
of the past.

To address this, we implemented exact leave-future-out cross-validation
(LFO-CV; Bürkner et al., 2020), adapted for our panel data structure. At
each step, the model was fit using data from blocks 1 through \emph{t},
and predictive performance was evaluated on block \emph{t}+1 for all
participants simultaneously. We set L=5 (25\% of blocks), ensuring
sufficient data for stable hierarchical parameter estimation --- a more
conservative proportion than the 12.5\% used in Bürkner et al.~(2020).

Across all three criteria --- WAIC, LOO, and LFO --- Mine1 consistently
emerged as the best-fitting model (Table X). Notably, the advantage of
Mine1 over Alt1 was considerably attenuated from LOO (ΔELPD = 99.8, SE =
12.4) to LFO (ΔELPD = 39.54, SE = 19.20). This attenuation is expected:
the log(trial) trend in Mine1 benefits disproportionately from future
information leakage under LOO, as future observations provide
information about the time trend that inflates apparent predictive
accuracy. The LFO estimate, free from such leakage, represents a more
honest assessment of the model's predictive advantage.

An additional finding concerns the comparison between Alt1 and the
Hypothesized model. The original study selected Alt1 over the
Hypothesized model based on WAIC, concluding that the separate modeling
of goal success and failure was warranted. Under LFO, however, this
advantage was not statistically distinguishable (ΔELPD = 38.57, SE =
33.80). This suggests that controlling for time trends may be more
consequential for model fit than distinguishing between success and
failure in the goal equation.

\subsection{The Emotion--Performance Link
Revisited}\label{the-emotionperformance-link-revisited}

The central finding of this reanalysis concerns the emotion →
performance pathway. Lin et al.~(2024) reported a significant negative
effect (β\_E→P = −0.103, 95\% HDI {[}−0.172, −0.036{]}) in their
best-fitting model (Alt1), which they interpreted through the lens of
the coasting hypothesis (Carver, 2003): positive emotions signal that
sufficient progress has been made, leading individuals to reduce effort
and consequently perform worse.

After introducing a log(trial) learning curve, this effect was
substantially attenuated and no longer credible (β\_E→P = −0.043, 95\%
CrI {[}−0.114, 0.028{]}). Simultaneously, a strong practice effect
emerged (β\_lt\_p = 0.326, 95\% CrI {[}0.262, 0.393{]}), indicating that
performance improved substantially over the course of the task. This
pattern suggests that the original negative emotion--performance
association may reflect a statistical confound rather than a genuine
causal process: as participants progressed through the repetitive math
task, performance increased due to practice while emotional engagement
declined --- two independent time trends that, when uncontrolled,
produce a spurious negative correlation between emotions and
performance.

Importantly, the remaining parameters of the model were largely robust
to this correction. The success/failure asymmetry in goal setting (β\_S
= 0.273, β\_F = 0.491) and the effect of the goal--performance
discrepancy on emotions (β\_GP→E = 0.128) were virtually unchanged,
suggesting that the core self-regulation dynamics identified by Lin et
al.~are genuine. It is specifically the emotion → performance pathway
that is not robust to the inclusion of a time trend.

\subsection{Reconciling with Study 2}\label{reconciling-with-study-2}

The interpretation of our findings gains further nuance when considered
alongside Study 2 of the original paper, in which German medical
students studied for a high-stakes second state exam over a period of up
to 40 days using an online learning platform (Breitwieser et al., 2022).
In contrast to the repetitive math task of Study 1 --- where
participants solved the same type of three-digit arithmetic problems
across 20 two-minute blocks in a single session --- Study 2 involved
varied learning content spread over weeks, with each study session
covering different medical topics. Furthermore, the task carried
substantial personal relevance, as performance was directly tied to
participants' professional futures.

Notably, Lin et al.~found that the direction of the emotion →
performance effect reversed in Study 2: when using postlearning
emotions, positive emotions predicted \emph{higher} subsequent
performance --- the opposite of Study 1. The authors acknowledged this
discrepancy but did not offer a unified explanation. Our findings
suggest one: in Study 1's repetitive task, a strong and monotonic
practice effect creates a powerful confound between time-trending
performance (upward) and time-trending emotions (downward, plausibly due
to fatigue or habituation to the task). The model, lacking a time trend
term, attributes this spurious covariation to β\_E→P.

In Study 2, by contrast, several features of the design would attenuate
this confound. The varied content across sessions makes systematic
practice-driven improvement less likely; the high stakes of the exam
sustain motivation and emotional engagement over time, reducing the
monotonic decline in affect seen in Study 1; and the temporal spacing of
sessions over days and weeks --- rather than minutes --- disrupts the
continuous accumulation of practice effects. Under these conditions, the
time-trend confound is weaker, and the observed emotion--performance
relationship may more closely reflect genuine affective processes.

This interpretation generates a testable prediction: introducing a time
trend into the Study 2 models should produce a smaller change in β\_E→P
than observed in Study 1. Unfortunately, we were unable to test this
prediction directly, as the Study 2 data were collected by Breitwieser
et al.~(2022) and were not publicly available at the time of this
analysis. We encourage future work to examine this hypothesis.

\subsection{Limitations}\label{limitations}

Several limitations should be noted. First, our model introduces the
learning curve as population-level (fixed) log(trial) effects rather
than allowing the time trend to vary across individuals. This choice was
made to maintain parsimony and avoid overparameterization: in
preliminary analyses, a random-slope specification for the log(trial)
effects encountered convergence difficulties, likely due to the limited
number of time points per individual (19 blocks) relative to the
additional individual-level parameters required. The fixed-effect
specification was sufficient to demonstrate the confound and represents
a conservative test --- if anything, allowing individual variation in
the learning curve would likely absorb additional variance currently
attributed to other parameters.

Second, we acknowledge that exact LFO-CV, while providing an unbiased
assessment of predictive performance, is computationally demanding.
Bürkner et al.~(2020) proposed an efficient PSIS approximation that
avoids refitting the model at most prediction steps by using
Pareto-smoothed importance sampling, with refitting only when the
approximation quality degrades (Pareto k \textgreater{} 0.7). However,
their framework was developed for univariate time series models. In our
case, the joint autoregressive structure across goals, performance, and
emotions --- combined with the hierarchical coupling of 357 participants
through shared population parameters --- means that each prediction step
introduces a large block of informative observations (357 × 3 outcome
variables). This renders the PSIS approximation infeasible, as the
posterior shift at each step is too large for importance sampling to
remain reliable. We therefore employed exact refitting at all 15
prediction steps. While tractable for our application (approximately 8
hours of computation per model), this approach would become prohibitive
for models with substantially longer time series or more expensive
individual fits. Developing efficient approximate LFO methods for
complex hierarchical autoregressive models remains an important
direction for future methodological work.

\subsection{Conclusion}\label{conclusion}

Our reanalysis supports the core theoretical framework proposed by Lin
et al.~(2024): goal setting, performance, and emotions are dynamically
interrelated, with the goal--performance discrepancy playing a central
role. However, the specific claim that positive emotions exert a
negative causal influence on subsequent performance --- interpreted via
the coasting hypothesis --- is not robust to the inclusion of a time
trend controlling for practice effects. Rather than reflecting a genuine
motivational process, this association appears to be an artifact of two
independent temporal trends: improving performance and declining
emotional engagement over the course of a repetitive task.

This finding carries practical implications for the design and analysis
of longitudinal self-regulation studies. Researchers employing
repeated-measures paradigms in which both performance and affect are
tracked over time should consider whether unmodeled time trends may
generate spurious associations between these variables. A simple
log(trial) covariate, as demonstrated here, can serve as a diagnostic
tool to assess the robustness of observed effects.

From a methodological standpoint, our work highlights the inadequacy of
WAIC and LOO-CV for model comparison in autoregressive panel models and
demonstrates that leave-future-out cross-validation, adapted for
hierarchical structures, provides a more principled assessment. We
encourage researchers working with longitudinal self-regulation models
to adopt similar diagnostic and model comparison practices. \#\#
Limitations and Future Directions

\appendix

\section{Model Specifications}\label{apx-a}

This appendix provides the full mathematical specification of the three
models compared in this study. All models share the same hierarchical
Bayesian structure and differ only in the specification of the
conditional means for goals, performance, and emotions. We first
describe the common elements, then present the model-specific equations.

\subsection{Common Structure}\label{common-structure}

At each time point \(t\) (blocks 2 through 20), goals (\(G\)),
performance (\(P\)), and emotions (\(E\)) are modeled as jointly
normally distributed, conditional on the previous time point:

\[
G_{t+1} \sim \mathcal{N}(\mu^G_{t+1},\; \sigma_G), \quad
P_{t+1} \sim \mathcal{N}(\mu^P_{t+1},\; \sigma_P), \quad
E_{t+1} \sim \mathcal{N}(\mu^E_{t+1},\; \sigma_E)
\tag{A1}
\]

where \(\sigma_G\), \(\sigma_P\), and \(\sigma_E\) are population-level
residual standard deviations. Block 1 values serve as initial conditions
and are not modeled. Each parameter \(\theta\) entering the conditional
means below is specified hierarchically via a non-centered
parameterization: \(\theta_i = \bar{\theta} + \sigma_\theta \cdot z_i\),
where \(z_i \sim \mathcal{N}(0, 1)\),
\(\bar{\theta} \sim \mathcal{N}(0, 3)\), and
\(\sigma_\theta \sim \mathcal{N}^+(0, 3)\). Subscript \(i\) indexes
participants.

\subsection{Conditional Means}\label{conditional-means}

The performance and emotion equations are shared across the Hypothesized
and Alt1 models:

\[
\mu^P_{t+1} = \gamma^P_i + \alpha^P_i \, P_t + \beta^{GP \to P}_i \,(G_{t+1} - P_t) + \beta^{E \to P}_i \, E_t \tag{A2}
\]

\[
\mu^E_{t+1} = \gamma^E_i + \alpha^E_i \, E_t + \beta^{GP \to E}_i \,(P_{t+1} - G_{t+1}) \tag{A3}
\]

The three models differ in the goal equation and in whether a log-trial
trend is included, as described below.

\textbf{Hypothesized Model.} The goal equation uses a single symmetric
coefficient for the goal-performance discrepancy (GPD):

\[
\mu^G_{t+1} = \gamma^G_i + \alpha^G_i \, G_t + \beta^{GP \to G}_i \,(P_t - G_t) + \beta^{E \to G}_i \, E_t \tag{A4}
\]

\textbf{Alt1 Model.} The goal equation replaces the symmetric GPD term
with separate coefficients for success and failure:

\[
\mu^G_{t+1} = \gamma^G_i + \alpha^G_i \, G_t + \beta^{E \to G}_i \, E_t +
\begin{cases}
\beta^{S}_i \,(P_t - G_t) & \text{if } P_t > G_t \text{ (success)} \\
\beta^{F}_i \,(P_t - G_t) & \text{if } P_t \leq G_t \text{ (failure)}
\end{cases}
\tag{A5}
\]

\textbf{LogTrial Model.} This model extends Alt1 by adding
population-level log-trial effects to all three conditional means. The
goal equation is identical to Equation A5 with an additional term
\(\beta^{\log t}_G \log(t+1)\). The performance and emotion equations
become:

\[
\mu^P_{t+1} = \gamma^P_i + \alpha^P_i \, P_t + \beta^{GP \to P}_i \,(G_{t+1} - P_t) + \beta^{E \to P}_i \, E_t + \beta^{\log t}_P \log(t+1) \tag{A6}
\]

\[
\mu^E_{t+1} = \gamma^E_i + \alpha^E_i \, E_t + \beta^{GP \to E}_i \,(P_{t+1} - G_{t+1}) + \beta^{\log t}_E \log(t+1) \tag{A7}
\]

The three log-trial coefficients (\(\beta^{\log t}_G\),
\(\beta^{\log t}_P\), \(\beta^{\log t}_E\)) are population-level only,
with no individual-level variation. This choice was made to maintain
parsimony; see Limitations for discussion.

\section{Another Appendix}\label{apx-b}






\end{document}
